\clearpage
\VerdierPhysicalPageAnchor{0261}%
\begingroup
\small
\setlength{\parindent}{0pt}%
\setlength{\parskip}{0pt}%
\setlength{\emergencystretch}{2em}%
\newcommand{\VerdierTermRow}[3]{%
  \VerdierSemanticAnchor{#1}%
  \par\noindent\hangindent=1.5em\hangafter=1
  #2。\enspace\textup{#3。}\par\smallskip}

\VerdierTermRow{idxt-p261-r01}{$n$-重複複体の二つの射を結ぶホモトピー}{第 I 章、2.5.1}
\VerdierTermRow{idxt-p261-r02}{$\operatorname{tr}_1(C)$ の二つの射を結ぶホモトピー}{第 I 章、3.1.4}
\VerdierTermRow{idxt-p261-r03}{ホモトピー同値}{第 I 章、2.5.11}
\VerdierTermRow{idxt-p261-r04}{遷移同型}{第 I 章、1.2.2}
\VerdierTermRow{idxt-p261-r05}{対象の $M$-次元}{第 III 章、2.2.5}
\VerdierTermRow{idxt-p261-r06}{$n$-重複複体の射}{第 I 章、2.1.4}
\VerdierTermRow{idxt-p261-r07}{次数 $g$ の射（$G$-圏における）}{第 I 章、1.3.8}
\VerdierTermRow{idxt-p261-r08}{次数付けと両立する次数付き関手の射}{第 I 章、1.1.1}
\VerdierTermRow{idxt-p261-r09}{$G$ の作用と両立する $G$-関手の射}{第 I 章、1.2.8}
\VerdierTermRow{idxt-p261-r10}{完全列に付随する連結射}{第 III 章、1.3.2}
\VerdierTermRow{idxt-p261-r11}{連結関手列の射}{第 III 章、3.3.3}
\VerdierTermRow{idxt-p261-r12}{三角の射}{第 I 章、3.2.1}
\VerdierTermRow{idxt-p261-r13}{分解可能な射}{第 II 章、1.2.8}
\VerdierTermRow{idxt-p261-r14}{スペクトル対象の射}{第 II 章、4.1.6}
\VerdierTermRow{idxt-p261-r15}{与えられた高さの初等射}{第 III 章、2.1.3}
\VerdierTermRow{idxt-p261-r16}{$\widetilde{\mathbb Z}$ の有限射}{第 II 章、4.3.1}
\VerdierTermRow{idxt-p261-r17}{$\mathcal{F}\ell_1(\mathcal D)$ の無視可能射}{第 II 章、3.1.1}
\VerdierTermRow{idxt-p261-r18}{分裂可能な射}{第 II 章、1.2.8}
\VerdierTermRow{idxt-p261-r19}{$n$-重複複体のホモトピックな射}{第 I 章、2.5.5}
\VerdierTermRow{idxt-p261-r20}{$\operatorname{tr}_1(C)$ のホモトピックな射}{第 I 章、3.1.4}
\VerdierTermRow{idxt-p261-r21}{完全関手またはコホモロジー的関手の核}{第 II 章、2.2.10}
\VerdierTermRow{idxt-p261-r22}{$\mathcal D^\wedge$ の許容対象}{第 II 章、3.2.2}
\VerdierTermRow{idxt-p261-r23}{$M$-次元が $n$ 以下の対象}{第 III 章、2.2.5}
\VerdierTermRow{idxt-p261-r24}{$F$-右（または左）自由化可能な対象}{第 II 章、2.3.2}
\VerdierTermRow{idxt-p261-r25}{$F$-右（または左）自由な対象}{第 II 章、2.3.2}
\VerdierTermRow{idxt-p261-r26}{別の対象上で $F$-右（または左）自由な対象}{第 II 章、2.3.2}
\endgroup
